% Image credits for the photographs and AI illustrations of Book 4 (University, Year 2), Hindi edition.
% Machine-readable license records live in images/book4/CREDITS.md; keep this
% file in sync with frontmatter/image-credits-book4.tex (one \item per item there).
% Names of people, institutions, collections, works, place names and licences are
% kept exactly as in the English credits (Latin script); only the frame and the
% descriptions are translated.
\chapter*{चित्रों का श्रेय}

इस पुस्तक के रेखाचित्र, मानचित्र और चार्ट मौलिक हैं; प्रकाशित आँकड़ों से बनाए गए मानचित्र और
चार्ट अपना स्रोत चित्र-परिचय में या इस सूची में बताते हैं। फ़ोटो उनके अपने-अपने मुक्त लाइसेंसों के
अंतर्गत, उनके लेखकों के प्रति आभार के साथ, इस्तेमाल किए गए हैं; हर फ़ोटो के पूरे स्रोत-लिंक और
लाइसेंस के पते पुस्तक के भंडार में \texttt{images/book4/CREDITS.md} में दिए गए हैं।

\bigskip
फ़ोटो और मौलिक रेखाचित्रों के साथ-साथ, कुछ चित्र AI से बनाए गए चित्रण हैं, जो पुस्तक के
संपादकों के लिखे निर्देशों से OpenAI के चित्र-निर्माण से बनाए गए, और शामिल करने से पहले उनकी
वैज्ञानिक शुद्धता जाँची गई। वे केवल सामान्य दृश्य दिखाते हैं, कभी कोई असली जगह, असली व्यक्ति
या असली नमूना नहीं। हर बनाए गए चित्र का पूरा निर्देश भंडार में
\texttt{images/book4/ai/PROMPTS.md} में दर्ज है।

\bigskip
\noindent फ़ोटो और आँकड़े, अध्याय के अनुसार:
\begin{itemize}\small\setlength{\emergencystretch}{2em}
  \item अध्याय~1 (अंतरतारकीय माध्यम और तारों का निर्माण): HL Tauri की आद्यग्रहीय चक्रिका, ALMA (ESO/NAOJ/NRAO), CC BY 4.0; Eagle नीहारिका में Pillars of Creation, NASA, ESA and the Hubble Heritage Team (STScI/AURA), सार्वजनिक डोमेन।
  \item अध्याय~2 (ग्रहों का निर्माण और उल्कापिंड): Allende कार्बनमय कॉन्ड्राइट और Gibeon लौह उल्कापिंड का विडमनश्टेटन प्रतिरूप, James St.~John, CC BY 2.0।
  \item अध्याय~3 (पृथ्वी में तत्त्व): बेसाल्ट में प्रावार का लर्ज़ोलाइट ज़ेनोलिथ, Kilbourne Hole, James St.~John, CC BY 2.0।
  \item अध्याय~5 (मैग्मा और उनका स्थापन): Shiprock के निकट एक डाइक, Dexter Perkins, CC0 (Wikimedia Commons)।
  \item अध्याय~6 (ज्वालामुखी उद्गार): Kīlauea के Pu`u~`\=O`\=o शंकु पर लावा का फ़व्वारा, J.\,D. Griggs, USGS, सार्वजनिक डोमेन; 18~मई 1980 को Mount St.~Helens का उद्गार-स्तंभ, Austin Post, USGS, सार्वजनिक डोमेन; 12~जून 1991 को Mount Pinatubo का उद्गार-स्तंभ, Dave Harlow, USGS, सार्वजनिक डोमेन।
  \item अध्याय~7 (प्रावार का पिघलना): ड्यूनाइट का हस्त-नमूना, Mike Davis (GeoDIL), CC0।
  \item अध्याय~8 (मैग्मीय विभेदन): Bushveld Complex की क्रोमाइटाइट परतें, Kevin Walsh, CC BY 2.0; Skaergaard अंतर्वेधन के स्तरित गैब्रो, C. Tegner, CC BY 4.0।
  \item अध्याय~10 (कायांतरी अभिक्रियाएँ और संलक्षणियाँ): लॉसनाइट नीलशिस्ट, Ward Creek, California, और एक्लोजाइट, Almenningen, Norway, James St.~John, CC BY 2.0।
  \item अध्याय~11 (तापदाबमिति और P--T--t पथ): गार्नेट शिस्ट, James St.~John, CC BY 2.0; एक्लोजाइट में कोसाइट का कण, आड़े ध्रुवकों के बीच, J.\,R. Smyth, CC BY-SA 1.0।
  \item अध्याय~12 (भूपर्पटी का गलन: मिग्मेटाइट और ग्रेनाइट): Nashua, New Hampshire में मिग्मेटाइट का अनावरण, Kevin M. Gill, CC BY 4.0; ल्यूकोग्रेनाइट शिरा, Anne Burgess (Geograph), CC BY-SA 2.0।
  \item अध्याय~13 (एक पर्वत-पट्टी की शारीरिकी: आल्प्स): Chenaillet ओफ़ियोलाइट के तकिया लावा, Lionel Duchoiselle, CC BY-SA 3.0; Zermatt से Matterhorn, Ximonic (Simo Räsänen), CC BY-SA 4.0।
  \item अध्याय~14 (टक्कर-पट्टियाँ और एंडीज़-प्रकार की पट्टियाँ: एक तुलना): Nepal और Shigatse के बीच तिब्बती पठार, Jan Kranendonk, CC BY-SA 4.0; Bolivia का Altiplano, Luca Galuzzi, CC BY-SA 2.5।
  \item अध्याय~15 (प्राचीन पुंजक और विल्सन चक्र): Siccar Point, दक्षिण-पूर्वी चट्टानी किनारा, dave souza, CC BY-SA 4.0।
  \item अध्याय~16 (अवसादी द्रोणियों की तुलना): North Sea में एक अर्ध-निमज्जनीय वेधन-प्लेटफ़ॉर्म, Erik Christensen, CC BY-SA 3.0।
  \item अध्याय~17 (महासागरीय द्वीप): International Space Station से Kure प्रवालद्वीप, NASA, सार्वजनिक डोमेन।
  \item अध्याय~18 (प्रीकैम्ब्रियन पृथ्वी): Pyke Hill, Ontario में स्पिनिफ़ेक्स कोमाटिआइट, James St.~John, CC BY 2.0; Nuccaleena Formation, South Australia का कैप कार्बोनेट, James St.~John, CC BY 2.0।
  \item अध्याय~19 (फ़ैनेरोज़ोइक पृथ्वी): Monument Rocks, Kansas के खड़िया-स्तंभ, James St.~John, CC BY 2.0; Bonarelli स्तर, Picenus, CC BY-SA 4.0।
  \item अध्याय~20 (जीवाश्म और पुरापर्यावरण): Burgess Shale से \emph{Waptia} और \emph{Marrella}, James St.~John, CC BY 2.0; \emph{Archaeopteryx} का Berlin नमूना, H.~Raab, CC BY-SA 3.0; Sagres में \emph{Zoophycos}, Amada44, CC BY 4.0।
  \item अध्याय~21 (वायुमंडलीय ऊष्मागतिकी): एक सुपरसेल तड़ित-झंझा का निहाई-शीर्ष, Stephen Corfidi, NOAA, सार्वजनिक डोमेन; Integrated Global Radiosonde Archive की रेडियोसोंड आरोहण-मापें, I. Durre और अन्य, NOAA NCEI, सार्वजनिक डोमेन।
  \item अध्याय~22 (मौसम प्रणालियाँ): Iceland के ऊपर एक निम्न-दाब तंत्र, NASA/GSFC, MODIS Rapid Response Team, Jacques Descloitres, सार्वजनिक डोमेन; International Space Station से Hurricane Isabel की आँख, Ed Lu, NASA, सार्वजनिक डोमेन।
  \item अध्याय~23 (जलवायु परिवर्तिता): उत्तर अटलांटिक दोलन और महासागरीय नीनो सूचकांक, NOAA Climate Prediction Center, सार्वजनिक डोमेन।
  \item अध्याय~24 (जलवायु अभिलेखागार और प्रॉक्सी): 1855~m पर GISP2 हिम-क्रोड की वार्षिक परतें, Anthony Gow, U.S. Army Corps of Engineers CRREL (NOAA Paleoclimatology), सार्वजनिक डोमेन; Diamond Caverns, Kentucky के एक ट्रैवर्टीन गुहा-निक्षेप की अनुप्रस्थ काट, James St.~John, CC BY 2.0; L.\,E. Lisiecki और M.\,E. Raymo (2005) का LR04 बेंथिक स्टैक, J. Jouzel और अन्य (2007) का EPICA Dome C तापमान-अभिलेख और B. Bereiter और अन्य (2015) का अंटार्कटिक \ce{CO2} संयुक्त अभिलेख, NOAA NCEI Paleoclimatology, सार्वजनिक डोमेन।
  \item अध्याय~25 (कक्षीय प्रणोदन और हिमयुग): Berger और Loutre (1991) के कक्षीय अवयव और सूर्यातप, NOAA NCEI Paleoclimatology, सार्वजनिक डोमेन।
  \item अध्याय~26 (भूवैज्ञानिक समय में जलवायु): सीनोज़ोइक बेंथिक समस्थानिक अभिलेख CENOGRID, T. Westerhold और अन्य (2020), PANGAEA, CC BY 4.0; G.\,L. Foster, D.\,L. Royer और D.\,J. Lunt (2017) का फ़ैनेरोज़ोइक \ce{CO2} संकलन, \emph{Nature Communications}, CC BY 4.0।
  \item अध्याय~27 (मानवजनित जलवायु परिवर्तन): IPCC Sixth Assessment Report, Working Group~I, Annex~III extended data, CC BY 4.0।
  \item अध्याय~28 (भूवैज्ञानिक संकट और जोखिम): Caraballeda, Venezuela में मलबा-प्रवाह के निक्षेप, Lawson Smith, U.S. Army Corps of Engineers, सार्वजनिक डोमेन; Aceh में 2004 की सुनामी के बाद खड़ी एक मस्जिद, USGS, सार्वजनिक डोमेन।
  \item अध्याय~29 (अयस्क निक्षेप): International Space Station से Bingham Canyon Mine, NASA (Expedition~15 crew), सार्वजनिक डोमेन; Flambeau Mine, Wisconsin का गॉसन, James St.~John, CC BY 2.0; D.\,A. Singer, V.\,I. Berger और B.\,C. Moring (2008) के ग्रेड--टनभार आँकड़े, USGS, सार्वजनिक डोमेन।
  \item अध्याय~30 (ऊर्जा भूविज्ञान): Geyserville, California के निकट एक भूतापीय बिजलीघर, U.S. Department of Defense, सार्वजनिक डोमेन; NIST Chemistry WebBook (Standard Reference Database 69) से कार्बन डाइऑक्साइड के घनत्व।
  \item अध्याय~31 (क्षेत्र और प्रयोगशाला विधियाँ II): प्लैजियोक्लेज़ को घेरे हुए पॉइकिलिटिक ऑगाइट, Strekeisen, CC BY-SA 4.0; एक गार्नेट--अभ्रक शिस्ट की पतली काट, Jackdann88, CC BY 3.0; Geological Survey of Finland में इलेक्ट्रॉन माइक्रोप्रोब, Kallerna, CC BY-SA 4.0; शैल संदर्भ पदार्थों के सिलिका परिणाम, J.\,E. Taggart और D.\,F. Siems (2002), USGS, सार्वजनिक डोमेन।
\end{itemize}
